近年，強化学習の実運用では環境変化に対し
短時間で性能を回復する適応性が求められる．
本研究では，方策再利用（Policy Reuse）
におけるエピソードごとの方策選択を多腕バンディットとして捉え，
従来のBoltzmann選択に加えてε-greedyとUCBを導入し，
格子状迷路で比較した．さらに，ゴール到達報酬のみと，
ステップ負報酬＋壁衝突ペナルティの2種の報酬設計で影響を検証した．
実験の結果，いずれの報酬設計でもPRQ系列は通常のQ学習を一貫して
上回り学習効率が向上した一方，最良の選択手法はタスクや
報酬構造に依存して変化することを示した．
以上より，方策選択の設計が再利用効果と安定性を左右することが示唆された．
